\chapter{Literature}\label{ch:lit}
This chapter reviews the literature on which the book draws. Section~\ref{I:sec:lit} reviews the literature of the whole series; the later sections add what is specific to Chapters~\ref{ch:II}, \ref{ch:III} and~\ref{ch:W}. Chapter~\ref{ch:IV} has no separate related-work section; its sources are discussed in its own text.

\begingroup
% macros of Part~I, Draft~11
\providecommand{\ov}{}\renewcommand{\ov}[1]{\overline{#1}}
\providecommand{\eps}{}\renewcommand{\eps}{\epsilon}
\providecommand{\A}{}\renewcommand{\A}{\mathbb{A}}
\providecommand{\C}{}\renewcommand{\C}{\mathbb{C}}
\providecommand{\Q}{}\renewcommand{\Q}{\mathbb{Q}}
\providecommand{\Z}{}\renewcommand{\Z}{\mathbb{Z}}
\providecommand{\cO}{}\renewcommand{\cO}{\mathcal{O}}
\providecommand{\gen}{}\renewcommand{\gen}{\mathrm{gen}}
\providecommand{\reg}{}\renewcommand{\reg}{\mathrm{reg}}
\providecommand{\tr}{}\renewcommand{\tr}{\operatorname{tr}}
\providecommand{\Gal}{}\renewcommand{\Gal}{\operatorname{Gal}}
\providecommand{\Aut}{}\renewcommand{\Aut}{\operatorname{Aut}}
\providecommand{\Nrd}{}\renewcommand{\Nrd}{\operatorname{Nrd}}
\providecommand{\Ad}{}\renewcommand{\Ad}{\operatorname{Ad}}
\providecommand{\Hom}{}\renewcommand{\Hom}{\operatorname{Hom}}
\providecommand{\Iso}{}\renewcommand{\Iso}{\operatorname{Iso}}
\providecommand{\vol}{}\renewcommand{\vol}{\operatorname{vol}}
\providecommand{\supp}{}\renewcommand{\supp}{\operatorname{supp}}
\providecommand{\diag}{}\renewcommand{\diag}{\operatorname{diag}}
\providecommand{\GL}{}\renewcommand{\GL}{\mathrm{GL}}
\providecommand{\SL}{}\renewcommand{\SL}{\mathrm{SL}}
\providecommand{\Gb}{}\renewcommand{\Gb}{\ov{G}}
\providecommand{\Zb}{}\renewcommand{\Zb}{\ov{Z}}
\providecommand{\GDb}{}\renewcommand{\GDb}{\ov{G}_D}
\providecommand{\pib}{}\renewcommand{\pib}{\ov{\pi}}
\providecommand{\ks}{}\renewcommand{\ks}{\kappa^{*}}
\providecommand{\Eb}{}\renewcommand{\Eb}{\ov{E}}
\providecommand{\Tb}{}\renewcommand{\Tb}{\ov{T}}
\providecommand{\cG}{}\renewcommand{\cG}{\mathcal{G}}
\providecommand{\cGb}{}\renewcommand{\cGb}{\ov{\mathcal{G}}}
\providecommand{\iD}{}\renewcommand{\iD}{\iota_D}
\providecommand{\iG}{}\renewcommand{\iG}{\iota_{\GL}}
\providecommand{\jj}{}\renewcommand{\jj}{\jmath}
\providecommand{\Pv}{}\renewcommand{\Pv}{\textbf{(P)}}
\providecommand{\Ev}{}\renewcommand{\Ev}{\textbf{(E)}}
\providecommand{\Ov}{}\renewcommand{\Ov}{\textbf{(O)}}
\providecommand{\Cv}{}\renewcommand{\Cv}{\textbf{(C)}}
% =====================================================================
\section{The literature of the series}\label{I:sec:lit}
This section reviews the literature on which the whole series draws, and places the series within it. Chapters~\ref{ch:II}--\ref{ch:IV} refer to this review and add only what is specific to them. It proves nothing and introduces no notation.

\paragraph{The Jacquet--Langlands correspondence.} Jacquet and Langlands \cite{JL70} related automorphic representations of the multiplicative group of a quaternion algebra to those of $\GL(2)$ by comparing trace formulas. For a central division algebra of arbitrary degree the local correspondence is due to Deligne--Kazhdan--Vign\'eras \cite{DKV84} and Rogawski \cite{Rog83}, and Badulescu \cite{Bad08} proved the global correspondence for all inner forms of $\GL(n)$, with multiplicity one and strong multiplicity one. Flicker and Kazhdan \cite{FK88} gave a simple trace formula for test functions with a single supercuspidal component and used it to extend both the simple-algebras correspondence and the metaplectic correspondence to automorphic forms with one supercuspidal component; Proposition~\ref{I:prop:FKvariant} is a statement of the same kind for the cover $\Gb(\A_K)$, proved here by following \cite[\S18]{FK86}, and we cite \cite{FK88} for its approach only. The transfer of orbital integrals between $\GL(n)$ and its inner forms, with explicit matching functions, is studied by Cohen \cite{Coh18}; harmonic-analytic background is in Kottwitz \cite{Kot05} and Harish-Chandra \cite{HC70}. Our definition of matching and the use of a simple trace formula follow the pattern of these works, with two new features: only good elements carry orbital integrals, and the two covers must first be identified on tori (\S\ref{I:sec:tori}). The sign convention is compared with Badulescu's in \S\ref{I:subsec:defmatch}.

\paragraph{Covering groups.} The $n$-fold covers of $\SL(2)$ and $\GL(2)$ over local and global fields go back to Kubota \cite{Kub67}, \cite{Kub69}. Moore \cite{Moo68} and Matsumoto \cite{Mat69} classified the topological central extensions of the split simply connected groups by Steinberg symbols (see also Milnor \cite{Mil71}); this classification is the main input of Proposition~\ref{I:prop:restriction}. Kazhdan and Patterson \cite{KP84} constructed the $n$-fold covers of $\GL(r)$ with a determinantal twisting parameter, developed their representation theory, including the exceptional representations, and formulated the generalized Shimura correspondence \cite{KP86}. Explicit block-compatible cocycles are in Banks--Levy--Sepanski \cite{BLS99}. The modern framework is that of Brylinski and Deligne \cite{BD01}, for which Weissman \cite{Wei18} and Gan--Gao \cite{GG18} formulated an $L$-group and a Langlands program; the historical introduction \cite{GGW18} surveys the subject. This series uses only the Kazhdan--Patterson covers and their restrictions to subgroups.

\paragraph{Metaplectic correspondences.} For $n=2$ the correspondence goes back to Shimura's work on modular forms of half-integral weight \cite{Shi73}; its representation-theoretic form is due to Gelbart \cite{Gel76} and Waldspurger \cite{Wal80}, \cite{Wal91}, through the Weil representation. Flicker \cite{Fli80} treated the $n$-fold cover of $\GL(2)$ by the trace formula, and Flicker and Kazhdan \cite{FK86} established the metaplectic correspondence between genuine representations of the $n$-fold cover of $\GL(r)$ and representations of $\GL(r)$, using a simple trace formula \cite[\S18]{FK86}; our $\GL$-side trace formula is theirs. Mezo generalized the local correspondence of \cite{FK86} and prepared a comparison of the invariant trace formulas of $\GL(r)$ and its covers \cite{Mez02}, \cite{Mez01}, and constructed metaplectic tensor products for irreducible representations \cite{Mez04}. Zou \cite{Zo23} revisited the local metaplectic correspondence, restored and generalized parts of the local theory of \cite{FK86}, and pointed out an error in \cite[\S26.2]{FK86}; Chapter~\ref{ch:III} uses his corrected form. Savin \cite{Sav88} proved a local Shimura correspondence at the level of Iwahori--Hecke algebras. W.-W.~Li developed the Arthur--Selberg trace formula for covers of connected reductive groups \cite{Li10}, \cite{Li12}, \cite{Li14}; it is one of the two routes to removing the auxiliary supercuspidal place (goal~(iii) of \S\ref{I:subsec:partII}). Kazhdan's lifting from tori \cite{Kaz84} has a metaplectic analog \cite{Jo88}, corrected in \cite{JC26a}, which is the source of the cover of $D^\times_{\A}$ used here. This chapter does not compose with the Flicker--Kazhdan correspondence; that would be the step after (D2).

\paragraph{Central extensions of division algebras.} The topological central extensions of $\SL_1(D)$, and the metaplectic kernel for anisotropic groups, were determined by Prasad--Raghunathan \cite{PR88} and Prasad--Rapinchuk \cite{PR96}; see also Rapinchuk \cite{Rap86}, \cite{Rap88} and Sury \cite{Sur05}. These works concern the extensions themselves rather than automorphic forms on them.

The work closest to ours on the side of the division algebra is that of Klose, in his G\"ottingen dissertation (1983, written under S.~J.~Patterson) and the two papers drawn from it. He studied central extensions by $\mu_n$ of the multiplicative group of a quaternion division algebra. The local paper \cite{Klo86a} concerns a $p$-adic quaternion division algebra; according to its indexing data, its subject is the commutator structure of $D^\times$, the structure of its group of one-units, and the classes in the second cohomology that describe its central extensions, including the metaplectic ones. The global paper \cite{Klo86b} concerns a quaternion division algebra $D$ over a number field. It constructs explicit cocycles, and in the description of Prasad and Rapinchuk \cite[Introduction]{PR96} it obtains partial results on the topological central extensions of the adele group of $\mathrm{GL}_{1,D}$; their own computation of the metaplectic kernel is more precise. Hill \cite[\S\S2.3, 3.8]{Hil08} attributes to Klose the facts on commutators and symmetric cocycles of abelian groups that were used there to determine the metaplectic extensions of $D^\times$; among them is a canonical isomorphism $H^2_{\mathrm{sym}}(\mu_m,\mu_m)\cong\Z/m$. In \cite{Jo88} it is noted that a result of Klose allows one to choose the quaternion algebra $D$ so that the cover of $G_D(\A_K)$ exists without assuming $(n,r)=1$, that is, also for even $n$. In this series $n$ is odd, so for $r=2$ that hypothesis holds automatically, and the cover of $D^\times$ is defined by restriction from the Kazhdan--Patterson cover of $\GL(r)$ over the splitting field $L$ (\cite{JC26a}, \S\ref{I:sec:covers}). We have not been able to consult \cite{Klo86a}, \cite{Klo86b} or the dissertation. The summary above rests on their indexing data (keywords and subject classification), on the descriptions in \cite{PR96}, \cite{Hil08} and \cite{Jo88}, and on the titles; the theorem statements of Klose's papers are not reproduced here for that reason. We have not compared Klose's extensions with the cover of \cite{JC26a} for $r=2$. Patel \cite{Pat14} shows that the $2$-fold metaplectic cover of $\GL(2,E)$, for a quadratic extension $E/F$, splits over $\GL(2,F)$ and over the multiplicative group of the quaternion division algebra over $F$; this is the same mechanism of restricting a cover over an extension field to a subgroup, and it is the mechanism by which our ambient cover realizes both $\GDb$ and $\Gb$.

\paragraph{The present series.} This chapter supplies the first steps of \cite[Remark~10.3]{JC26a}. We found no existing metaplectic Jacquet--Langlands correspondence for covers of inner forms of $\GL(r)$. In particular, among the papers that zbMATH lists as citing \cite{FK86} (Zbl 0616.10024), only \cite{Jo88} concerns a cover of $D^\times$; those mentioning division or simple algebras treat linear groups. Chapter~\ref{ch:II} proves the matching of arbitrary functions, using germ expansions \cite[\S8]{FK86} and results on the linear groups (Shalika germs, Steinberg pseudo-coefficients), and improves Theorem~C. Chapter~\ref{ch:III} studies the division-algebra side through a unitary transform of class functions, using Adams operations at tame places and Kottwitz's framework for rational conjugacy. Chapter~\ref{ch:IV} reconciles the cocycle conventions with \cite{FK86} and \cite{Zo23}, studies power maps and toral positivity at the places dividing $n$, and proves a conditional global comparison. The companion chapter (Chapter~\ref{ch:W}) studies the remaining wild local problem in degree two, using the Heisenberg--Weil formalism and character theory of compact $p$-adic groups.
\endgroup
\begingroup
% macros of Part~II, Draft~10
\providecommand{\ov}{}\renewcommand{\ov}[1]{\overline{#1}}
\providecommand{\A}{}\renewcommand{\A}{\mathbb{A}}
\providecommand{\C}{}\renewcommand{\C}{\mathbb{C}}
\providecommand{\Z}{}\renewcommand{\Z}{\mathbb{Z}}
\providecommand{\cO}{}\renewcommand{\cO}{\mathcal{O}}
\providecommand{\fp}{}\renewcommand{\fp}{\mathfrak{p}}
\providecommand{\gen}{}\renewcommand{\gen}{\mathrm{gen}}
\providecommand{\reg}{}\renewcommand{\reg}{\mathrm{reg}}
\providecommand{\tr}{}\renewcommand{\tr}{\operatorname{tr}}
\providecommand{\Nrd}{}\renewcommand{\Nrd}{\operatorname{Nrd}}
\providecommand{\Ad}{}\renewcommand{\Ad}{\operatorname{Ad}}
\providecommand{\Iso}{}\renewcommand{\Iso}{\operatorname{Iso}}
\providecommand{\vol}{}\renewcommand{\vol}{\operatorname{vol}}
\providecommand{\supp}{}\renewcommand{\supp}{\operatorname{supp}}
\providecommand{\Cent}{}\renewcommand{\Cent}{\operatorname{Cent}}
\providecommand{\GL}{}\renewcommand{\GL}{\mathrm{GL}}
\providecommand{\St}{}\renewcommand{\St}{\mathrm{St}}
\providecommand{\Gb}{}\renewcommand{\Gb}{\ov{G}}
\providecommand{\Zb}{}\renewcommand{\Zb}{\ov{Z}}
\providecommand{\GDb}{}\renewcommand{\GDb}{\ov{G}_D}
\providecommand{\pib}{}\renewcommand{\pib}{\ov{\pi}}
\providecommand{\Eb}{}\renewcommand{\Eb}{\ov{E}}
\providecommand{\Tb}{}\renewcommand{\Tb}{\ov{T}}
\providecommand{\cG}{}\renewcommand{\cG}{\mathcal{G}}
\providecommand{\cGb}{}\renewcommand{\cGb}{\ov{\mathcal{G}}}
\providecommand{\iD}{}\renewcommand{\iD}{\iota_D}
\providecommand{\iG}{}\renewcommand{\iG}{\iota_{\GL}}
\providecommand{\jj}{}\renewcommand{\jj}{\jmath}
\providecommand{\sg}{}\renewcommand{\sg}{\sigma}
\providecommand{\Pv}{}\renewcommand{\Pv}{\textbf{(P)}}
\providecommand{\Ev}{}\renewcommand{\Ev}{\textbf{(E)}}
\providecommand{\Ov}{}\renewcommand{\Ov}{\textbf{(O)}}
\providecommand{\Cv}{}\renewcommand{\Cv}{\textbf{(C)}}
% =====================================================================
\section{Work related to \texorpdfstring{Chapter~\ref{ch:II}}{Chapter 5}}\label{II:sec:lit}
This section discusses the work related to Chapter~\ref{ch:II}. The literature of the whole book (the Jacquet--Langlands correspondence, covering groups, metaplectic correspondences, central extensions of division algebras) is reviewed in \S\ref{I:sec:lit}; we add only what is specific to Chapter~\ref{ch:II}. The section proves nothing and introduces no notation.

The local results of Chapter~\ref{ch:II} concern harmonic analysis near the center, and their inputs are of two kinds. For covers of $\GL(r)$, Kazhdan--Patterson \cite[\S1]{KP86} and Flicker--Kazhdan \cite[\S8]{FK86} proved the germ expansion and the characterization of orbital integrals of genuine functions (after Vign\'eras), and \cite{FK86} constructed the local metaplectic correspondence for square-integrable representations; Zou \cite{Zo23} corrected and extended parts of that theory. W.-W.~Li developed the local harmonic analysis needed for the trace formula on covers \cite{Li12}, including semisimple descent, the Plancherel formula and the density of tempered characters, and proved the invariant trace formula for covers \cite{Li14}. On the linear side we use the Shalika germ expansion \cite{Sha72}, Howe's work on the Fourier transform and germs of characters for $\GL_n$ \cite{How74}, Harish-Chandra's theory of admissible invariant distributions \cite{HC70}, \cite{HC78}, Rogawski's computation of the trivial germ on elliptic tori via the building \cite{Rog81}, van Dijk's formula for induced characters \cite{vD72}, and, for parabolic descent of germs, Kivinen--Tsai \cite{KT26}. Pseudo-coefficients of square-integrable representations and the density theorem are taken in the form stated in \cite[\S\S25.1, 27.3]{FK86}, which goes back to Kazhdan \cite{Kaz86}; the linear Jacquet--Langlands inputs are from \cite{DKV84}, \cite{Rog83}, \cite{Bad02} and \cite[Theorem~2.2]{Bad08}. The Plancherel formula and Harish-Chandra's finiteness theorem are used in the form of Waldspurger \cite{Wal03}, with explicit formulas for $\GL(n,F)$ and $\GL(m,D)$ in Aubert--Plymen \cite{AP05}; and Kottwitz's normalization of measures on inner forms is in \cite{Kot88}. For the classification of tempered representations of $\GL(r,F)$ we use Zelevinsky \cite{Zel80}; general background on smooth representations is in Bernstein--Zelevinsky \cite{BZ77} and Bernstein \cite{Ber84}.

Our matching theorem uses none of the deeper parts of the theory of covers. Its main tool is the section $\sg$ of \S\ref{II:sec:center}, which reduces genuine harmonic analysis near the center of either cover to linear harmonic analysis, compatibly with the identification of elliptic tori of Chapter~\ref{ch:I}. That covers split over small open subgroups is standard \cite[\S2]{FK86}; what is used here is the stronger statement that the splitting given by $n$-th powers of $n$-th roots is equivariant and consistent with $\jj$ on a full conjugation-invariant neighborhood of the unipotent set, at every finite place including those dividing $n$.

% =====================================================================
\endgroup
\begingroup
% macros of Part~III, Draft~10
\providecommand{\A}{}\renewcommand{\A}{\mathbb A}
\providecommand{\C}{}\renewcommand{\C}{\mathbb C}
\providecommand{\Z}{}\renewcommand{\Z}{\mathbb Z}
\providecommand{\GL}{}\renewcommand{\GL}{\operatorname{GL}}
\providecommand{\Nrd}{}\renewcommand{\Nrd}{\operatorname{Nrd}}
\providecommand{\Irr}{}\renewcommand{\Irr}{\operatorname{Irr}}
\providecommand{\vol}{}\renewcommand{\vol}{\operatorname{vol}}
\providecommand{\tH}{}\renewcommand{\tH}{\widetilde H}
\providecommand{\tZ}{}\renewcommand{\tZ}{\widetilde Z}
\providecommand{\ttT}{}\renewcommand{\ttT}{\widetilde T}
\providecommand{\reg}{}\renewcommand{\reg}{\mathrm{reg}}
\providecommand{\gen}{}\renewcommand{\gen}{\mathrm{gen}}
\providecommand{\JL}{}\renewcommand{\JL}{\mathrm{JL}}
\providecommand{\Sha}{}\renewcommand{\Sha}{\operatorname{Sha}}
\providecommand{\Gb}{}\renewcommand{\Gb}{\overline{G}}
\providecommand{\GDb}{}\renewcommand{\GDb}{\overline{G}_D}
\providecommand{\Zb}{}\renewcommand{\Zb}{\overline{Z}}
\providecommand{\Ptag}{}\renewcommand{\Ptag}{\textnormal{\textbf{(P)}}}
\providecommand{\Etag}{}\renewcommand{\Etag}{\textnormal{\textbf{(E)}}}
\providecommand{\Otag}{}\renewcommand{\Otag}{\textnormal{\textbf{(O)}}}
\providecommand{\Ctag}{}\renewcommand{\Ctag}{\textnormal{\textbf{(C)}}}
\section{Work related to \texorpdfstring{Chapter~\ref{ch:III}}{Chapter 6}}\label{III:sec:related}
This section discusses the work related to Chapter~\ref{ch:III}. The literature of the whole book is reviewed in \S\ref{I:sec:lit}; we add only what is specific to Chapter~\ref{ch:III}. The section proves nothing and introduces no notation.

\paragraph{The division-algebra side.} The local Jacquet--Langlands correspondence for $D^\times$ is due to Deligne, Kazhdan and Vign\'eras \cite{DKV84}, and Badulescu \cite{Bad08} proved the global correspondence for inner forms of $\GL_n$; Cohen \cite{Coh18} studies the transfer of representations and of orbital integrals between $\GL_n$ and its inner forms. The norm-one group $H^1=\mathrm{SL}_1(D)$ of a $p$-adic division algebra was studied by Riehm \cite{Rie70}, and its topological central extensions by Prasad and Raghunathan \cite{PR88}; Klose \cite{Klo86a} studied commutators and the second cohomology of $p$-adic quaternion division algebras, including their metaplectic extensions (see \S\ref{I:sec:lit}). Proposition~\ref{III:prop:tame-split} (unique splitting over $H^1$ when $p\nmid n$) and Proposition~\ref{III:prop:wild-obstruction} (non-splitting when $p\mid n$) are statements of this kind for the particular cover of Chapter~\ref{ch:I}; they are proved here directly, and we cite these works for their subject only.

\paragraph{Power operations.} The Adams operations come from topological $K$-theory \cite{Ada62}; Atiyah's account of power operations \cite{Ati66} and Segal's equivariant $K$-theory \cite{Seg68} give the representation-ring form used in \S\ref{III:sec:adams} and in the tame proof (Lemma~\ref{III:lem:tame-adams}).

\paragraph{Metaplectic correspondences.} The square-integrable metaplectic correspondence for $\GL_r$ is due to Flicker and Kazhdan \cite{FK86}; Zou \cite{Zo23} restored and generalized it and gave the root-sum relation used in Corollary~\ref{III:cor:D2tame}, whose constant is compared with $b_D$ in Remark~\ref{III:rem:normalizations}. The transfer $L_D$ is defined through the $n$-th power map on good elements and the central character, after the model of Kazhdan's lifting and its metaplectic analog (\S\ref{I:sec:lit}).

\paragraph{Arithmetic.} Section~\ref{III:sec:kottwitz} follows Kottwitz \cite{Kot84} on rational and stable conjugacy, transfer of tori to inner forms and Tamagawa factors; the Galois cohomology of $\Sha^1$ is in Milne \cite{Mil06}. Proposition~\ref{III:prop:global-power-hasse} is a statement modulo scalars of the kind given, for $n$-th powers themselves, by the classical local--global principle whose sharp form is the Grunwald--Wang theorem \cite[Ch.~X]{AT68}; for odd $n$ the exceptional case of that theorem does not arise. The $p$-adic change-of-variables formula used in \S\ref{III:sec:tame} is quoted from \cite{Loe08}.
\endgroup
\begingroup
% macros of the companion chapter on the rank-two wild place, Draft~8
\providecommand{\F}{}\renewcommand{\F}{F}
\providecommand{\Z}{}\renewcommand{\Z}{\mathbb Z}
\providecommand{\C}{}\renewcommand{\C}{\mathbb C}
\providecommand{\Q}{}\renewcommand{\Q}{\mathbb Q}
\providecommand{\Fq}{}\renewcommand{\Fq}{\mathbb F}
\providecommand{\Nm}{}\renewcommand{\Nm}{\operatorname N}
\providecommand{\Nrd}{}\renewcommand{\Nrd}{\operatorname{Nrd}}
\providecommand{\Trd}{}\renewcommand{\Trd}{\operatorname{Trd}}
\providecommand{\Tr}{}\renewcommand{\Tr}{\operatorname{Tr}}
\providecommand{\Irr}{}\renewcommand{\Irr}{\operatorname{Irr}}
\providecommand{\Ind}{}\renewcommand{\Ind}{\operatorname{Ind}}
\providecommand{\Ad}{}\renewcommand{\Ad}{\operatorname{Ad}}
\providecommand{\tH}{}\renewcommand{\tH}{\widetilde H}
\providecommand{\tZ}{}\renewcommand{\tZ}{\widetilde Z}
\providecommand{\tU}{}\renewcommand{\tU}{\widetilde U}
\providecommand{\tV}{}\renewcommand{\tV}{\widetilde V}
\providecommand{\tth}{}\renewcommand{\tth}{\widetilde\theta}
\providecommand{\Ocal}{}\renewcommand{\Ocal}{\mathcal O}
\providecommand{\Pfr}{}\renewcommand{\Pfr}{\mathfrak P}
\providecommand{\pfr}{}\renewcommand{\pfr}{\mathfrak p}
\providecommand{\Reg}{}\renewcommand{\Reg}{\operatorname{Reg}}
\providecommand{\Heis}{}\renewcommand{\Heis}{\mathrm{Heis}}
\providecommand{\Cubes}{}\renewcommand{\Cubes}{\mathrm{Cubes}}
\providecommand{\Good}{}\renewcommand{\Good}{\mathrm{Good}}
\providecommand{\Ptag}{}\renewcommand{\Ptag}{\textnormal{\textbf{(P)}}}
\providecommand{\Pctag}{}\renewcommand{\Pctag}{\textnormal{\textbf{(P$^{\mathrm c}$)}}}
\providecommand{\Etag}{}\renewcommand{\Etag}{\textnormal{\textbf{(E)}}}
\providecommand{\Otag}{}\renewcommand{\Otag}{\textnormal{\textbf{(O)}}}
\providecommand{\Ctag}{}\renewcommand{\Ctag}{\textnormal{\textbf{(C)}}}
\section{Work related to \texorpdfstring{Chapter~\ref{ch:W}}{Chapter 8}}\label{W:sec:literature}
This section places Chapter~\ref{ch:W} among the works it draws on or runs
parallel to.  It is selective; \S\ref{I:sec:lit} has the fuller
discussion of the metaplectic correspondence itself.

\paragraph{Metaplectic covers.}  The $n$-fold covers of $\mathrm{SL}_2$
and $\mathrm{GL}_2$ over a local field containing $\mu_n$ go back to
Kubota \cite{Kub67}; \cite{Kub69}, who gave the explicit cocycle used in
Section~\ref{W:sec:model}; Moore \cite{Moo68} and Matsumoto classified the
topological central extensions of simply connected groups by Steinberg
symbols, which is why the Hilbert symbols of $F$ and $E_0$ could be
pinned down in Lemma~\ref{W:lem:model} as the unique Steinberg symbols.
Kazhdan and Patterson \cite{KP84} constructed the $n$-fold covers of
$\mathrm{GL}_r$ with the determinantal twist parameter $c$, developed
their representation theory (the metaplectic tensor product, the
exceptional representations), and formulated the generalized Shimura
correspondence; the parameter $c_v=-1/4\bmod p$ of
Section~\ref{W:sec:targetC}, forced by scalar centrality, is a point on
that family.  The modern framework is that of Brylinski--Deligne
central extensions \cite{BD01}, for which Weissman \cite{Wei18} and Gan--Gao
\cite{GG18} formulated an $L$-group and a Langlands program; a Satake
isomorphism for covers is in \cite{McN12}; \cite{Gao18}.  The cover of $D^\times$ used
here is defined by restriction from $\mathrm{GL}_2(E_0)$ (Chapter~\ref{ch:I}); we are
not aware of an independent literature on $n$-fold covers of division
algebras beyond \cite{JC26a} and this series.

\paragraph{Metaplectic correspondences.}  For $n=2$ the Shimura
correspondence was realized by Waldspurger \cite{Wal80}; \cite{Wal91} through the
theta correspondence, as recalled in Section~\ref{W:subsec:howe}; Flicker
\cite{Fli80} treated the $n$-fold cover of $\mathrm{GL}_2$ by the trace
formula, and Flicker--Kazhdan \cite{FK86} established the metaplectic
correspondence between the $n$-fold cover of $\mathrm{GL}_r$ and
$\mathrm{GL}_r$ in general, at the cost of a simple trace formula and
hypotheses at the places dividing $n$.  Kazhdan's lifting \cite{Kaz84} and
its metaplectic analogue \cite{JC26a} are the local models for the division-side
transfer $L_D$ of Chapter~\ref{ch:III}: a transfer of characters defined on good
elements through the $n$-th power map and the central character.  At
the level of the Iwahori--Hecke algebra, Savin \cite{Sav88} proved a local
Shimura correspondence for the tame covers; Zou \cite{Zo23} revisits the
local metaplectic correspondence and its applications.  None of these
works addresses integrality at a place dividing $n$ for a division
algebra, which is the problem of Targets~A--C; Chapter~\ref{ch:III}'s tame argument
is an Adams-operation argument in the sense of Section~\ref{W:subsec:howe},
and the obstruction at $p\mid n$ is the one described in
Remark~\ref{W:rem:dependencies} and Section~\ref{W:sec:future}.

\paragraph{Representations of division algebras and Jacquet--Langlands.}
The irreducible representations of $D^\times$ were constructed by Corwin
\cite{Cor74} and Howe \cite{How77b} in the tame case and by Koch--Zink \cite{KZ80} and
Corwin in general; the local Jacquet--Langlands correspondence is due to
Jacquet--Langlands \cite{JL70} for $r=2$ and to Deligne--Kazhdan--Vign\'eras
\cite{DKV84} and Rogawski \cite{Rog83} in general, with the global correspondence
completed by Badulescu \cite{Bad08}.  Bushnell and Henniart made the
correspondence explicit in the essentially tame case \cite{BH11} and treat
$\mathrm{GL}_2$ and the quaternion algebra in detail in \cite{BH06}, including
the conductor and depth bookkeeping behind the notion of level in
Section~\ref{W:subsec:divalg}.  The induced models of
Sections~\ref{W:sec:maintheorem} and~\ref{W:sec:plan} are of the shape of
Howe's construction, with a character of a covered congruence subgroup
in place of a character of a torus.

\paragraph{Heisenberg groups, Weil representations and the Howe
correspondence.}  Weil \cite{Wei64} introduced the metaplectic
representation; Howe \cite{How79} formulated the duality conjecture, proved
by Waldspurger \cite{Wal90} for $p\ne2$ and by Gan--Takeda \cite{GT16} in general;
\cite{MVW87} is the standard reference for the $p$-adic theory.  The finite
Weil representations and their characters are in G\'erardin \cite{Ger77}.
We use this circle of ideas only as a guide (Section~\ref{W:subsec:howe}
and Remark~\ref{W:rem:howe-duality}); the Stone--von Neumann theorem for
covered tori is used literally (\cref{I:prop:good}).

\paragraph{Characters of compact $p$-adic groups.}  Howe's Kirillov
theory \cite{How77a} describes the characters of congruence subgroups by
coadjoint orbits when the group is large relative to $p$; the orbit
method for uniform pro-$p$ groups with $p$ odd was developed further by
Jaikin-Zapirain \cite{Jai06} and Boyarchenko--Sabitova \cite{BS08}, and the
cohomology of uniform groups by Lazard \cite{Laz65} (see [DDMS; SW]), which
enters the alternative route to (H1) in Remark~\ref{W:rem:H1-checks}.  The
finite-group tools are Brauer's characterization of characters and
Clifford theory \cite{Isa76}.  The Hilbert-symbol facts of
Section~\ref{W:sec:model} are in Fesenko--Vostokov \cite{FV02} and Serre \cite{Ser79}.

\paragraph{Computation.}  The exact computations behind the \Pctag\
labels are small-scale: each is a short, purpose-written Python
script -- for instance an orbit enumeration on a specific group of
order $3^9$, or a character table obtained numerically and then
verified exactly against the class-sum eigenvector criterion --
rather than a call to a general-purpose computer-algebra library
(Appendix~\ref{W:app:scripts}).  A general-purpose system such as
GAP or Magma could recompute the
character tables of Section~\ref{W:sec:plan} directly, and the design
search of Section~\ref{W:sec:finite} is an enumeration that any system can
reproduce.  The methodological choice to separate hand proofs,
exhaustive exact computations and sampled corroboration, and to label
them, follows the practice of the series.
\endgroup
